\section{Nachgereichter Text B: Korrekturen und ein neuer prüfbarer
Feldadapter}\label{nachgereichter-text-b-korrekturen-und-ein-neuer-pruxfcfbarer-feldadapter}

Der Text mit 45 Zeilen wurde vollständig gelesen. Seine unveränderte
Kopie steht in \codeword{late\_sources/attachment\_B.txt}, SHA-256
\codeword{9aaf5e18cbc6a84d610b88429872df9b69e6078365bfdc8133c8f17e2c1c3801}.
Die für die neue Prüfung maßgebliche \codeword{field\_dictionary.py} ist
gesondert eingefroren, SHA-256
\codeword{a581bacd17c08dd9a766ae439bfe1e8e35e6d30b9509cb9d98053b897d8ecc13}.
Die bereits abgeschlossenen Schatten-Dateien wurden nicht verändert.

\subsection{Übernahmeentscheidung}\label{uxfcbernahmeentscheidung}

{\def\LTcaptype{none} % do not increment counter
\begin{longtable}[]{@{}
  >{\raggedright\arraybackslash}p{(\linewidth - 2\tabcolsep) * \real{0.5000}}
  >{\raggedright\arraybackslash}p{(\linewidth - 2\tabcolsep) * \real{0.5000}}@{}}
\toprule\noalign{}
\begin{minipage}[b]{\linewidth}\raggedright
Aussage im Text
\end{minipage} & \begin{minipage}[b]{\linewidth}\raggedright
Belastbare Fassung / konkrete Folge
\end{minipage} \\
\midrule\noalign{}
\endhead
\bottomrule\noalign{}
\endlastfoot
1,44 Milliarden unsichtbare Freiheiten im kleinsten nichttrivialen
Sektor & 1444233216 ist die komplexe Kommutantdimension des gewährten
Kontrollsatzes X,Nb im N=3-Sektor. Im N=2-Sektor lautet die
entsprechende Zahl 3829536. Weder Zahl zählt direkt physische
Freiheitsgrade. \\
Mit voller Symmetrie sieben, mit Modenzählern eins; damit nichts
Physisches verborgen & Die assoziativen Algebren sind unter den
ausdrücklich ergänzten Ressourcen bestimmt. Ein skalarer Kommutant
beweist weder vollständige dynamische Lie-Kontrolle noch ausführbare
Präparation/Messinstrumente. Ein neuer exakter Spin-1-Gegenzeuge unten
trennt diese Aussagen. \\
Die Kontrollfrage ist eine einzige Ja/Nein-Frage & Bereits v1.6.6/7
enthalten Zwischenstufen: einzelner Dunkelprojektor, ein getrennter
Casimir, nur SU(4), nur Spin(10). Auch X und Nb unabhängig schalten zu
können ist stärker als ein fixes H. \\
Nur beide getrennten Casimire können die drei dunklen Typen
unterscheiden & Bereits einer der getrennten Casimire genügt, da seine
drei Eigenwerte verschieden sind. Der Gesamtcasimir allein genügt
nicht. \\
Der Skalarkanal ist endgültig verboten; Vermittler muss (1,0) sein &
Exakt im lokalen, ableitungsfreien Ansatz gleicher Weylhändigkeit mit
unverändertem antisymmetrischem W und ohne zusätzliche unabhängige
Komponenten. Der symmetrische Kanal trägt dort. Dies ist kein
vollständiges Feldwörterbuch jedes möglichen Kontinuumsadapters. \\
Das Lochkomposit besitzt einen brauchbaren Spin-1/2-Projektor & Der
angegebene Projektor ist für den \textbf{gleichhändigen}
Tensorproduktraum korrekt. Für das wörtlich adjungierte Weylfeld ist
zunächst die Händigkeit zu klären; die neue exakte Boost-Prüfung unten
macht die Konsequenz sichtbar. \\
N=0 modulo vier verkürzt die Grundzustandskonkurrenz auf
Singulettsektoren & Die Regel gilt für Singuletts. Erst ein unabhängiger
globaler Eindeutigkeitssatz erzwingt für die zusammenhängende
halbeinfache Gruppe einen Singulett-Grundzustand. Ohne ihn kann ein
tieferes nichtsinguläres, entartetes Multiplett gewinnen. Zudem ist die
Liste bosonischer Ladungssektoren nach oben unbeschränkt. \\
480, 916, 299520/229 sind drei exakte Energiesprossen & Es sind
dimensionslose quadrierte Kopplungskoeffizienten, nach Ausklammern von
g², keine drei Eigenenergien. Sie stimmen mit dem Anfang der richtigen
Voll-H-Lanczoskette überein. \\
Nur die gesamte 33-Sprossen-Leiter fehlt noch & 33 zählt die
Bosonzahlstufen k=0,\ldots,32 bei N=64. Eine Zustandsrichtung je Stufe
schließt unter H nicht. Schon k=2 hat vier Singuletts; das
Rücklaufresiduum besitzt exakt positive Norm 5001523200/229. Die
tatsächliche vierte Voll-H-Lanczosdiagonale ist 105168998/26292551·Δ,
nicht 4Δ. \\
Ritzüberlappung 39,7 Prozent bestätigt den wahren Grundzustand & Der
Wert gehört zum benannten Ritz-Zustand. Ohne Residuum und
Spektralabstand ist er keine entsprechende Überlappungsaussage über Ω.
Die zitierte Ritzenergie setzt außerdem g/Δ=1/20 voraus. \\
Schwächere Nachbarsektorfloors lassen Grundzustandssatz wieder offen &
Sie lassen diese schwächere Sonde offen. Das widerlegt oder ersetzt den
stärkeren vorherigen Satz im unveränderten Modellvertrag nicht. \\
Operationssatz und Feldwörterbuch fertig; räumliche Skalierung jetzt
freigegeben & Die erforderlichen Ressourcen und der Lorentzadapter sind
weiterhin konkrete offene Verbindungen. Die neue gemeinsame
Schattenkarte bestimmt deren nächste Prüfung wesentlich genauer als
diese pauschale Freigabe. \\
\end{longtable}
}

Die normierte direkte Summe der drei Feldkanäle hat weiterhin Norm 8√3
bei Normquadrat 192. Die eingefrorene Feldquelle enthält noch die
bereits korrigierte Addition von drei Normen zu 24; ihr gesamtes grünes
Ergebnis wurde daher nicht als unkritischer Vertrauensbeweis übernommen.

\subsection{1. Neuer exakter
Kontrollgegenzeuge}\label{neuer-exakter-kontrollgegenzeuge}

Die gewöhnlichen drei Spin-1-Matrizen J\_x,J\_y,J\_z auf C³ erfüllen
{[}J\_x,J\_y{]}=iJ\_z und zyklisch. Bereits J\_x und J\_z haben
ausschließlich skalare gemeinsame Kommutanten. Der dynamische Lie-Raum
ist trotzdem nur su(2) mit Dimension drei, nicht su(3) mit Dimension
acht. Das wird im neuen Prüfer durch exakte Ränge und alle drei
Kommutatoren kontrolliert.

Damit ist die Unterscheidung praktisch: Ein assoziativer Abschluss kann
vollständig sein, während die tatsächlich durch Hamiltonpulse
erreichbaren Unitären eine echte Untergruppe bilden. Für den nativen
Kontrollvertrag braucht es weiterhin die konkret gewährten Generatoren,
ihre dynamische Lie-Algebra und eine ausgeführte
Puls-/Instrumentenfolge.

\subsection{2. Was am Spin-1/2-Projektor
stimmt}\label{was-am-spin-12-projektor-stimmt}

Die Feldquelle verwendet die sechs unnormierten Koordinaten

\codeword{(00,0), (00,1), (01,0), (01,1), (11,0), (11,1)},

wobei \codeword{(01)} die Summe \codeword{|01>+|10>} bezeichnet. Die
korrekte Gram-Matrix ist

\[\adjustbox{max width=\linewidth}{$\displaystyle 
 G_6=\operatorname{diag}(1,1,2,2,1,1).
$}\]

Der neue Prüfer baut den vollständigen Symmetrisierer auf drei
Spinorindizes unabhängig neu auf. Zugleich konstruiert er die direkte
Epsilon-Auslese

\[\adjustbox{max width=\linewidth}{$\displaystyle 
 C=\begin{pmatrix}0&1&-1&0&0&0\\0&0&0&1&-1&0\end{pmatrix},
 \qquad R=\frac23G_6^{-1}C^T.
$}\]

Exakt gelten

\[\adjustbox{max width=\linewidth}{$\displaystyle 
 CR=I_2,\quad P_{1/2}=RC,\quad
 P_{3/2}=I_6-P_{1/2},\quad
 CP_{3/2}=0,\quad CP_{1/2}=C.
$}\]

Beide Projektoren sind idempotent, komplementär und haben Rang zwei
beziehungsweise vier. P\_\{1/2\} ist selbstadjungiert bezüglich G\_6;
seine nicht symmetrische 6×6-Koordinatenmatrix ist kein Fehler. Die
Epsilon-Auslese und Wiedereinsetzung intertwinen exakt alle drei
sl(2)-Generatoren. Dies liefert eine ausführbare algebraische
Extraktionsregel im \textbf{angenommenen gleichhändigen} Raum

\[\adjustbox{max width=\linewidth}{$\displaystyle 
 (1,0)\otimes(1/2,0)=(3/2,0)\oplus(1/2,0).
$}\]

Die Zahl 1/3 ist nur der Ranganteil im isotropen Zustand. Der neue
Prüfer zeigt positive reine Vektoren mit Spin-1/2-Gewicht null und eins.
Deshalb ersetzt der Projektorrang keine Berechnung seines Gewichts im
nativen Grundzustand oder im geladenen Spektralpol.

\subsection{3. Die neue präzise
Händigkeitsschranke}\label{die-neue-pruxe4zise-huxe4ndigkeitsschranke}

Die eingefrorene Tabelle in \codeword{field\_dictionary.py}, Zeile 327,
bezeichnet den zweiten Faktor als
\codeword{f\textasciicircum{}dagger (1/2,0)}. Dafür muss die
Feldabbildung explizit festgelegt werden. Für ein tatsächliches
adjungiertes linkes Weylfeld gilt dagegen ψ† in (0,1/2). Hermitesche
Konjugation vertauscht diese Darstellungen; siehe
\href{https://arxiv.org/pdf/0812.1594\#page=9}{Dreiner, Haber und
Martin, Abschnitt 2}.

In diesem \textbf{wörtlichen Feldadjungierten-Fall} ist das Produkt

\[\adjustbox{max width=\linewidth}{$\displaystyle 
 (1,0)\otimes(0,1/2)=(1,1/2)
$}\]

ein irreduzibler Lorentzraum mit komplexer Dimension sechs. Unter der
Rotationsuntergruppe zerfällt er weiterhin in Spin 3/2 und Spin 1/2;
diese beiden Rotationsräume sind aber nicht getrennt boostinvariant. Der
Prüfer verifiziert den skalaren Kommutanten sämtlicher sechs Generatoren
des komplexifizierten Lorentzraums; die beiden unabhängigen
sl(2)-Faktoren wirken als irreduzibles äußeres Tensorprodukt.

Ein konkreter Zeuge genügt bereits. In obiger Basis setze

\[\adjustbox{max width=\linewidth}{$\displaystyle 
 J_z=J_z^{(1)}\otimes I_2+I_3\otimes J_z^{(1/2)},\qquad
 K_z=J_z^{(1)}\otimes I_2-I_3\otimes J_z^{(1/2)}.
$}\]

Bis auf den konventionellen Faktor i ist K\_z der Boostgenerator für
entgegengesetzte Händigkeit. Dann gilt exakt

\[\adjustbox{max width=\linewidth}{$\displaystyle 
 [P_{1/2},J_z]=0,\qquad
 [P_{1/2},K_z]=
 \begin{pmatrix}
 0&0&0&0&0&0\\
 0&0&4/3&0&0&0\\
 0&-2/3&0&0&0&0\\
 0&0&0&0&2/3&0\\
 0&0&0&-4/3&0&0\\
 0&0&0&0&0&0
 \end{pmatrix}.
$}\]

Dieser Kommutator hat Rang vier. Der Rotationsprojektor kann deshalb in
diesem Vertrag nicht als Lorentz-kovarianter Weyl-Ausleseprojektor
verwendet werden.

\textbf{Das ist kein pauschaler Ausschluss des ursprünglichen
Modenkomposits.} Ein bloßer Fock-Erzeuger f† ist noch kein adjungiertes
lokales Weylfeld. Bereits die übliche linke Weylfeldentwicklung enthält
Erzeugungs- und Vernichtungsterme mit passenden Spinorwellenfunktionen;
siehe \href{https://arxiv.org/pdf/0812.1594\#page=24}{Dreiner, Haber und
Martin, Gleichungen 3.1.3-3.1.5}. Eine unabhängige oder passend
ladungskonjugierte gleichhändige Komponente kann den gleichhändigen
Projektor erhalten. Sie muss aber als konkreter Adapter mit korrekter
Ladung, CAR, W-Kontraktion und Dynamik gezeigt werden. Aus dem Symbol †
allein folgt weder ihr Vorliegen noch ihr Ausschluss.

\subsection{4. Konkreter nächster
Annahmetest}\label{konkreter-nuxe4chster-annahmetest}

Der Anschlussauftrag ist jetzt endlich und entscheidbar formuliert:

\begin{enumerate}
\def\labelenumi{\arabic{enumi}.}
\tightlist
\item
  Den tatsächlichen Feldadapter für \codeword{D\_r ∼ b f†}
  beziehungsweise den davon verschiedenen \codeword{χ† ∼ b†ηf†}
  hinschreiben: gepunktete/ungepunktete Indizes, Ladung, Fourieranteile,
  erlaubte Zusatzkomponenten.
\item
  Für genau diesen Adapter alle Rotations- \textbf{und Boostgeneratoren}
  auf dem zusammengesetzten Träger berechnen. Ein Kandidat mit
  Lorentz-Auslese C muss \codeword{C L\_composite = L\_Weyl C} erfüllen.
  Der neue gleiche-/gegengängige Test liefert dafür positive und
  negative Referenzfälle.
\item
  Die Kontraktion am unveränderten W ausführen und anschließend das
  tatsächliche Matrixelement \codeword{C D\_r†|Ω>} beziehungsweise den
  entsprechenden Ladungskanal berechnen. Erst dessen normiertes
  Spektralgewicht kann den vorhandenen geladenen Pol mit einem
  Spin-1/2-Feld verbinden.
\item
  Falls der wörtlich adjungierte Fall gewählt wird, keinen bloßen
  Rotationsprojektor als Lorentzprojektor verwenden. Ein zusätzlicher
  Impuls-/Ableitungsadapter oder eine andere Händigkeit kann untersucht
  werden, zählt aber als ausdrückliche neue Struktur.
\end{enumerate}

Dieser Test ergänzt die gemeinsame Schattenrekonstruktion: Die dort
bewiesene vollständige innere Tomographie ersetzt die hier fehlende
Lorentz-Intertwining-Abbildung nicht. Umgekehrt definiert ein richtiger
Lorentzprojektor noch kein ausführbares natives Messinstrument.

\subsection{Reproduktion}\label{reproduktion}

\begin{verbatim}
python3 -B replay_late_field.py
\end{verbatim}

Benötigt werden nur Python und SymPy. Der separate Lauf verifiziert zwei
Quellenpins sowie \textbf{36 exakte Bedingungen pro Modus}, keine
numerischen Bedingungen. Normaler Lauf und \codeword{-OO} ergeben
byte-identische JSON-Ausgaben; Warnungen gelten als Fehler und alle
Guards bleiben unter Optimierung aktiv. Die vollständigen Matrizen
stehen in \codeword{late\_field.normal.json}.

Das Ergebnis-JSON hat SHA-256
\codeword{cfd9e15cec5c2b0ce3de9397614836b16b24b4adb8f6711eb704925589965774}.
Der Prüfer hat SHA-256
\codeword{f18e6fd2af7c1a28a81e85654df68fe805e59b1d957801611d2668678550dfc2}.
Die fremde Feldsuite wird nur eingefroren und punktuell unabhängig
geprüft, nicht als vollständig fehlerfreie Gesamtreproduktion
ausgegeben.

\begin{center}\rule{0.5\linewidth}{0.5pt}\end{center}

\section{Nachtragsprüfung C: fundamentale Reduktion und globale
Quelle}\label{nachtragspruxfcfung-c-fundamentale-reduktion-und-globale-quelle}

\section{Nachtrag C: überwiegend bereits abgedeckte fundamentale
Reduktion}\label{nachtrag-c-uxfcberwiegend-bereits-abgedeckte-fundamentale-reduktion}

\subsection{Übernahmeurteil}\label{uxfcbernahmeurteil}

Der neue Text ist eine ausführlichere Fassung der bereits eingefrorenen
fundamentalen Reduktion, ergänzt um die weitgehend gleiche frühere
Gesprächszusammenfassung. Er bringt \textbf{keinen neuen H-Term, keine
neue Operationsquelle und keinen neuen Transferbeweis}. Seine
wesentlichen mathematischen Einschränkungen sind richtig und bleiben in
der gemeinsamen Darstellung erhalten. Eine neue Großrechnung ist dafür
nicht notwendig.

Neue Eingabe: 299 Zeilen, SHA-256
\codeword{a10f702306c7409535df9f4d26f0995970833a62297a1cb93d8c95cdce746d4d}.
Quelle:
\codeword{/Users/stefanhamann/.codex/attachments/b0d7e533-3ca8-4a84-966c-2546b066dd08/pasted-text.txt}.

Bereits eingefrorene Kurzfassung:
\codeword{universalraum-singlet-observable-20260915/late\_sources/fundamental\_reduction.txt},
99 Zeilen, SHA-256
\codeword{0081cebd9759e8fa776995ee8efb3231c7611a82bea9cdeea86967073d1459a2}.
Die Dateien sind nicht bytegleich; der Neuheitsbefund betrifft ihren
Inhalt.

\subsection{Was bestätigt und bereits abgedeckt
ist}\label{was-bestuxe4tigt-und-bereits-abgedeckt-ist}

{\def\LTcaptype{none} % do not increment counter
\begin{longtable}[]{@{}
  >{\raggedright\arraybackslash}p{(\linewidth - 2\tabcolsep) * \real{0.5000}}
  >{\raggedright\arraybackslash}p{(\linewidth - 2\tabcolsep) * \real{0.5000}}@{}}
\toprule\noalign{}
\begin{minipage}[b]{\linewidth}\raggedright
Aussage
\end{minipage} & \begin{minipage}[b]{\linewidth}\raggedright
Urteil und Grenze
\end{minipage} \\
\midrule\noalign{}
\endhead
\bottomrule\noalign{}
\endlastfoot
Volle Schur-Antwort T†F(H−E0)T=γF I64 & Richtig bei gruppeninvariantem H
und Grundzustand sowie irreduziblem ursprünglichem Entnahmemultiplett.
Gilt für die volle beschränkte Spektralfunktion, nicht nur einen
dominanten Pol. Schon in v1.6.7 BIG\_PICTURE B2 und der eingefrorenen
Kurzfassung enthalten. \\
Aus C\_rs(t)=δ\_rs c(t) folgt kein Ortsindex in diesen 64 Labels &
Richtig. Eine andere lineare Benennung erzeugt nur einen Gramfaktor.
Keine pauschale Sperre gegen Vielteilchendynamik, nichtinvariante
Referenzzustände oder andere Operatorfamilien. \\
Symmetriekommutant und Operationskommutant haben verschiedene Bedeutung
& Richtig und bereits unterschieden. Erstgenannter erlaubt Dynamik
zwischen Vorkommen desselben Typs; ein großer Kommutant eines
eingeschränkten Operationssatzes kann fehlenden Zugriff anzeigen. Seine
Größe beweist weder Raum noch ausführbare Kontrolle. \\
μN ist bei festem Gesamt-N eine Konstante & Richtig. Auf demselben
präparierten Eingang bleiben freie Heisenbergentwicklungen neutraler
Observablen identisch. Für ganze Interventionsfolgen muss auch deren
zugelassener Gesamtvertrag respektiert werden. \\
Eine globale Zustandsregel braucht keinen äußeren Präparator & Richtig.
Das Verbot, mit N-erhaltenden Operationen aus N=0 nach N=64 zu gelangen,
ist ein Satz über einen bestimmten Präparationsweg. Es verbietet weder
eine globale Randbedingung noch eine andere begründete Zustandsregel.
Interne Detektoren und Präparationen bleiben trotzdem herzuleiten. \\
Keine nichttriviale Fermion-Teilparität in einer nativen Bank bei festen
Bosonen & Bereits exakt geprüft: zusammenhängender W-Paargraph und
binärer Rang 63. Ein Paritätsverbot getrennt kopierter Banken darf nicht
auf beliebige Teilmengen derselben Bank übertragen werden. \\
\end{longtable}
}

Das eigene v1.6.8-Ergebnis ergänzt diese Diagnose: Der bedingte lokale
Phasenimpuls und die terminale Tomografie benutzen einen gemeinsamen
invarianten N=2-Dreizustandsraum. Dessen Startzustand ist nicht der in
Schurs Aussage vorausgesetzte invariante N=64-Grundzustand. Es gibt
daher keinen Widerspruch und keine Berechtigung, die neue Umverteilung
auf dessen geladenen Einlochpol oder auf räumliche Bewegung zu
übertragen.

